% TOP_PROBLEM: {"id": 2660, "title": "Kirby Problem 1.1", "category_id": 7, "category": {"id": 7, "name": "topology", "display_name": "Topology", "description": "Properties preserved under continuous deformations.", "slug": "topology", "order_index": 7, "created_at": "2026-07-31T15:26:25.671Z"}, "status": "open", "problem_number": "KP-1.1", "set_id": 11}
% FIRST_POSTED: 2026-10-01
% ENTRY_KIND: statement_only
% UnsolvedMath ID 2660; source catalogue code KP-1.1
% !TeX program = lualatex
% Definition and sources only; this file is not an LLM solution attempt.
\documentclass[11pt,a4paper]{article}
\usepackage[margin=25mm]{geometry}
\usepackage{fontspec}
\setmainfont{lmroman10-regular.otf}[BoldFont=lmroman10-bold.otf,
 ItalicFont=lmroman10-italic.otf,BoldItalicFont=lmroman10-bolditalic.otf]
\usepackage{amsmath,amssymb,mathtools}
\usepackage{unicode-math}
\setmathfont{latinmodern-math.otf}
\AtBeginDocument{\let\setminus\smallsetminus}
\usepackage{xurl,hyperref}
\hypersetup{colorlinks=true,urlcolor=blue}
\setcounter{secnumdepth}{0}
\setlength{\parindent}{0pt}
\setlength{\parskip}{5pt}
\setlength{\emergencystretch}{3em}
\begin{document}
\section*{Kirby Problem 1.1}\label{2660}
{\small UnsolvedMath ID 2660 · KP-1.1 · Topology · Kirby's Problems in Low-Dimensional Topology}
\subsection{Definitions and mathematical statement}
A knot is a smooth embedding of a circle in the oriented sphere $S^3$, considered
up to ambient isotopy. A diagram is a generic projection to a plane, with an
overpassing strand specified at every transverse double point. If $D$ is such a
diagram, let $c(D)$ be its number of double points. The crossing number is
\[
c(K)=\min\{c(D):D\text{ is a diagram of }K\}.
\]
For oriented knots $K_1,K_2$, their connected sum $K_1\#K_2$ is obtained by deleting
a short unknotted arc from each and joining the resulting endpoints without
introducing additional knotting. Its oriented isotopy class is independent of
these choices.

The question is whether every pair of knots satisfies
\[
c(K_1\#K_2)=c(K_1)+c(K_2).
\]
Joining minimal diagrams proves the upper bound. The requested content is the
reverse inequality for an arbitrary diagram of the composite knot; that diagram
need not visibly display its summands. The assertion includes nonalternating knots
and arbitrary numbers of summands by induction.
\subsection{Short English statement}
Does tying two knots together always add their minimum crossing numbers?
\subsection{Sources}
\begin{itemize}
\item[S1] ULAM.AI and the UnsolvedMath Contributors, supplied problems.json, record 2660 (KP-1.1), Kirby's Problems in Low-Dimensional Topology. Catalogue identity and source transcription; CC BY 4.0.
\url{https://huggingface.co/datasets/ulamai/UnsolvedMath}.
\item[S2] K3: A New Problem List in Low-Dimensional Topology, Problem 1.1. Expanded formulation of the supplied catalogue statement.
\url{https://math.berkeley.edu/sites/default/files/surv-295-ruberman-watermarked-author-pdf.pdf}.
\end{itemize}
\end{document}
